% Reconstructed from the 14 August 2026 FINAL PDF.
\documentclass[11pt]{article}
\usepackage[margin=1.1in]{geometry}
\usepackage{amsmath,amssymb,amsthm}
\usepackage{enumitem}
\usepackage[colorlinks=true,linkcolor=blue,citecolor=blue,urlcolor=blue]{hyperref}

\newcommand{\pP}{\textbf{[P]}}
\newcommand{\pC}{\textbf{[C]}}
\newcommand{\pO}{\textbf{[O]}}

\title{\textbf{Status of the High-Precision Coefficient of the\\
Universal Kinematic Residue $I_B$\\
(Version 4)}}
\author{Robert W.\ McGwier\\
\small CTO and Staff Scientist, Cohere Technology Group, Herndon, VA}
\date{August 14, 2026 --- Version 4}

\begin{document}
\maketitle

\begin{abstract}
\noindent
This note records the final status of the high-precision numerical
coefficient of the universal kinematic residue $I_B$. Controlled
numerical experiments in the $d=2$ laboratory, including
singularity-subtracted and principal-value prescriptions, establish
that the relative residue is of order unity. A multi-digit coefficient
remains open: a fully analytic Hadamard expansion of the current
two-point function under modular flow is still required. The existence
of a non-zero residue of relative size $O(1)$ is \pP{}; the precise
digits are \pO{}. No structural claim of the Emergent Gravity Program
depends on those digits.

Why this note is required. Earlier versions left the coefficient
entirely open or reported only qualitative non-vanishing. The present
version documents the best controlled numerical evidence (order-unity
relative residue) and states with precision what remains to be done for
a high-precision extraction, so that the public record is complete and
honest.
\end{abstract}

\section*{Epistemic Conventions}
Every material claim is tagged \pP{}, \pC{} or \pO{}.
Numerical agreement is not proof.

\section{Controlled Numerical Evidence}
Experiments were performed with the exact M\"obius modular flow of an
interval and the chiral current two-point function. Smooth compactly
supported sources were used. Both direct subtraction of the local
quadratic form and principal-value / finite-part prescriptions for the
difference of flowed and unflowed correlators were examined.

Stable results under all variations:
\begin{itemize}[leftmargin=1.4em]
\item the residue is non-vanishing;
\item the integrand is even;
\item exponential decay at large modular time confirms infrared safety;
\item the relative residue (residue divided by the local quadratic form)
is of order unity ($O(1)$).
\end{itemize}
The best controlled estimate of the relative residue under the
singularity-subtracted prescriptions examined is approximately $1$--$5$,
depending on the precise finite-part implementation. This establishes
the order of magnitude but not a multi-digit coefficient.

\section{What Remains Open}
A fully controlled high-precision coefficient requires an analytic
Hadamard (or equivalent distributional) expansion of the current
two-point function under the modular flow, isolating and subtracting
all terms that produce non-integrable singularities when folded against
the modular kernel $\pi/(2\cosh^{2}(\pi s))$. The numerical proxies used
here cancel the leading $1/r^{2}$ piece but do not yet remove every
singular contribution systematically. Until that analytic step is
completed, the precise digits remain \pO{}.

\section{Logical Status}
\begin{itemize}[leftmargin=1.4em]
\item Existence of a non-zero residue of relative size $O(1)$: \pP{}.
\item Absorbability of the residue into a finite dictionary improvement:
\pP{} (established independently of the precise coefficient).
\item Multi-digit numerical coefficient of $I_B$: \pO{}.
\end{itemize}
The precise coefficient controls only the size of the dictionary
improvement. It is not required for the validity of the torsionful first
law, the unconditional metric sector, the sign of the Hehl--Datta
interaction, or the pure-information determination of its magnitude.
All of those claims remain certified independently of the digits.

\section*{Final Note on the Numerical Campaign}
Further purely numerical iteration with hard cutoffs or simple
principal-value proxies will not produce a reliable multi-digit
coefficient. The remaining task is analytic. This Version 4 therefore
closes the numerical campaign and records the order-of-magnitude result
together with the precise requirement for any future high-precision
extraction.

\begin{thebibliography}{99}
\bibitem{IV} R.W.\ McGwier, ``Condition NL for Finite Balls,'' Zenodo (2026).
\bibitem{VI} R.W.\ McGwier, ``Evaluation of the Universal Kinematic Integral\ldots,'' Zenodo (2026).
\bibitem{VIII} R.W.\ McGwier, ``Status of the High-Precision Coefficient\ldots'' (earlier versions), Zenodo (2026).
\bibitem{CTT} H.\ Casini, E.\ Teste and G.\ Torroba, J.\ Phys.\ A 50 (2017) 364001.
\bibitem{CHM} H.\ Casini, M.\ Huerta and R.C.\ Myers, JHEP 05 (2011) 036.
\bibitem{FGHMV} T.\ Faulkner et al., JHEP 03 (2014) 051.
\end{thebibliography}
\end{document}
